\section{第四章\quad 理想气体混合物及湿空气}
% 只考 1、2

若各组分气体均处于理想气体状态，则混合气体也可作为理想气体处理。混合气体可等效为某种“假想气体”，其总质量、总物质的量等于各组分之和。

\pt{混合气体状态方程} $pV=mR_{\mathrm{g,eq}}T$;
\pt{平均摩尔质量与折合气体常数关系}：$M_{\mathrm{eq}}R_{\mathrm{g,eq}}=R$；
\pt{物质的量与平均摩尔质量}：$n_{\mathrm{eq}}=\sum_i n_i,n_{\mathrm{eq}}M_{\mathrm{eq}}=\sum_i n_iM_i$ 因此：$M_{\mathrm{eq}}=\sum_i x_iM_i,R_{\mathrm{g,eq}}=\frac{R}{M_{\mathrm{eq}}}$，$R_{\mathrm{g,eq}}=\sum_i w_iR_{\mathrm{g},i},M_{\mathrm{eq}}=\frac{1}{\sum_i \frac{w_i}{M_i}}$

\pt{道尔顿分压定律}：$p=\sum_i p_i$；
\pt{分容积定律}：$V=\sum_i V_i$
\pt{混合气体的质量分数}：$w_i=\frac{m_i}{m},\qquad \sum_i w_i=1$ ；
\pt{混合气体的摩尔分数}：$x_i=\frac{n_i}{n},\qquad \sum_i x_i=1$ ；
\pt{混合气体的体积分数}：$y_i=\frac{V_i}{V},\qquad \sum_i y_i=1$ ；
\pt{混合气体的比热容}：$c_{\mathrm{mix}}=\sum_i w_ic_i$；
\pt{混合气体的热力学能}：$U_{\mathrm{mix}}=\sum_i U_i$；
\pt{混合气体的焓}：$H_{\mathrm{mix}}=\sum_iH_i=\sum_i(U_i+p_iV)=U_{\mathrm{mix}}+pV$；
\pt{混合气体熵}：$S_{\mathrm{mix}}=\sum_iS_i=\sum_i m_is_i$；
无化学反应混合气体过程熵变，按 $1\,\mathrm{kg}$ 混合气体：$\mathrm{d}s_{\mathrm{mix}}=\sum_i w_i\mathrm{d}s_i=\sum_i w_i\left(c_{p,i}\frac{\mathrm{d}T}{T}-R_{\mathrm{g},i}\frac{\mathrm{d}p_i}{p_i}\right)$；
定比热容时：$  \Delta s_{\mathrm{mix}}=\sum_i w_ic_{p,i}\ln\frac{T_2}{T_1}-\sum_i w_iR_{\mathrm{g},i}\ln\frac{p_{i,2}}{p_{i,1}}$；
按 $1\,\mathrm{mol}$ 混合气体：$\mathrm{d}S_{\mathrm{m,mix}}=\sum_i x_iC_{p,\mathrm{m},i}\frac{\mathrm{d}T}{T}-\sum_i x_iR\frac{\mathrm{d}p_i}{p_i}$

\divider

大气由 $\mathrm{N_2}$、$\mathrm{O_2}$、$\mathrm{Ar}$、$\mathrm{CO_2}$、水蒸气及微量气体组成；
若空气中水蒸气为过热状态，$t>t_s(p_{\mathrm{v}})$，则湿空气未饱和；若空气中水蒸气达到饱和，$t=t_s(p_{\mathrm{v}})$，则湿空气饱和；\pt{露点}：湿空气中水蒸气分压力 $p_{\mathrm{v}}$ 所对应的饱和温度称为露点温度

\divider
